% Comparação de duas funções afins com mesmo b e a diferentes (uma crescente, outra decrescente).
% Usado em: Matematica/9ano/unidade-4-funcao-afim/capitulo_02_coeficientes-e-comportamento.md (final da seção 2.2)
\begin{tikzpicture}[scale=0.7]
  % Eixos
  \draw[->,thick] (-4,0) -- (4,0) node[right,font=\small] {$x$};
  \draw[->,thick] (0,-1.5) -- (0,7) node[above,font=\small] {$y$};

  % Marcações
  \foreach \i in {-3,-2,-1,1,2,3}{
    \draw (\i,0.12) -- (\i,-0.12) node[below,font=\scriptsize] {\i};
  }
  \foreach \i in {1,2,3,4,5,6}{
    \draw (0.12,\i) -- (-0.12,\i) node[left,font=\scriptsize] {\i};
  }
  \node[below left,font=\scriptsize] at (0,0) {$O$};

  % Reta crescente: f(x) = 2x + 1 (a > 0)
  \draw[very thick,eleveBlue,domain=-2.5:3,smooth] plot (\x, {2*\x + 1});
  \node[font=\small,eleveBlue,fill=white,inner sep=1pt] at (3.0,7.4) {$f(x)=2x+1$};

  % Reta decrescente: g(x) = -3x + 1 (a < 0, mesmo b)
  \draw[very thick,eleveAccent,domain=-2:2,smooth] plot (\x, {-3*\x + 1});
  \node[font=\small,eleveAccent,fill=white,inner sep=1pt] at (-2.6,7.4) {$g(x)=-3x+1$};

  % Ponto comum (0, 1) destacado
  \fill[black] (0,1) circle (3.5pt);
  \node[font=\scriptsize,fill=white,inner sep=1pt] at (0.7,1.3) {$(0,\,1)$};

  % Setas indicando crescente/decrescente
  \draw[->,eleveBlue,thick] (1.6,5.0) to[bend left=10] (2.6,6.5);
  \node[font=\scriptsize,eleveBlue] at (2.4,4.5) {sobe};

  \draw[->,eleveAccent,thick] (-1.6,5.0) to[bend right=10] (-1.0,3.0);
  \node[font=\scriptsize,eleveAccent] at (-2.2,4.5) {desce};
\end{tikzpicture}
